\documentclass[12pt]{article}
\usepackage{a4wide}
\usepackage[T2A]{fontenc}
\usepackage[utf8]{inputenc}
\usepackage[english,russian]{babel}
\usepackage{amssymb,amsmath,bm}
\usepackage{graphicx}
\usepackage{tikz-cd}
\usetikzlibrary{babel}
\usepackage{enumitem}
\usepackage[font=small]{caption}
\usepackage{hyperref}
\graphicspath{{figures/}}

\newcommand{\vx}{\mathbf{x}}
\newcommand{\tht}{\bm{\theta}}

\begin{document}

\section*{Вычислительный эксперимент}

Цель: исследовать устойчивость алгоритмов восстановления структуры суперпозиции к шуму в матрице смежности. Шум вносится только в матрицу; значения целевой переменной в выборке точные. Эксперимент для одной модели состоит из следующих шагов:
\begin{enumerate}[nosep]
	\item По истинной формуле модели построить её матрицу смежности $M$.
	\item Зашумить матрицу: $N=M+\alpha U$, $U_{ij}\sim\mathcal U(-1,1)$, $\alpha\in[0,1]$, затем нормировать $P=(N-\min N)/(\max N-\min N)$.
	\item Восстановить матрицу $\hat M=R(P)$ одним из алгоритмов $R$.
	\item Проверить качество восстановления по каждому из критериев раздела <<\nameref{sec:criteria}>>. Для проверки по качеству описания данных:
	\begin{enumerate}
		\item Задать значения параметров $\tht$ и по истинной формуле построить выборку $D$ (интегрирование ОДУ либо снимок полей для МГД).
		\item Построить по $\hat M$ модель $f_{\hat M}$, подобрать её параметры $\hat\tht$ по выборке $D$.
		\item Сравнить $f_{M,\tht}(X)$ и $f_{\hat M,\hat\tht}(X)$.
	\end{enumerate}
\end{enumerate}

Модели в эксперименте "--- осциллятор Ван дер Поля, системы Лотки--Вольтерры и Лоренца и уравнения МГД. Для всех уровней шума $\alpha$ используется одна и та же реализация $U$.

Выборки для ОДУ получены интегрированием траекторий со случайными параметрами и начальными условиями (рис.~\ref{fig:phase} и \ref{fig:lorenz-phase}). Данные для МГД "--- гладкие синтетические поля на сетке $12^3$. Целевые значения вычисляются по правым частям уравнений. Для ОДУ "--- 20 выборок, 20 реализаций шума на выборку и 41 уровень шума. Для МГД "--- 10, 10 и 21 соответственно. Прогоны Лоренца учитывают исправление повторного выбора листьев в DFS и BFS. Для остальных моделей приведены результаты до исправления.

\begin{figure}[!ht]
\centering
\includegraphics[width=\linewidth]{vdp_phase}\\[2mm]
\includegraphics[width=\linewidth]{lv_phase}
\caption{Фазовые портреты и траектории. Сверху "--- осциллятор Ван дер Поля (хаотический режим), снизу "--- модель Лотки--Вольтерры (замкнутые циклы). Слева "--- номинальные значения параметров, справа "--- несколько выборок с параметрами из окрестности номинала}
\label{fig:phase}
\end{figure}

\begin{figure}[!htbp]
\centering
\includegraphics[width=\linewidth]{lorenz_phase}
\caption{Система Лоренца. Слева проекция фазового портрета на плоскость $(x,y)$ при номинальных параметрах, справа временные ряды $x(t)$ для трёх выборок параметров из априорных распределений}
\label{fig:lorenz-phase}
\end{figure}

Сравниваются DFS, BFS, алгоритм Прима и варианты с PCST (для осциллятора "--- только DFS, BFS и Прим). Точный словарь содержит только функции исходной модели, расширенный "--- ещё и лишние функции. На рис.~\ref{fig:trees} показаны примеры деревьев при $\alpha=0.7$. При таком шуме восстановленное дерево может быть правдоподобнее истинного по зашумлённой матрице.

\begin{figure}[!ht]
\centering
\includegraphics[width=0.49\linewidth]{vdp_tree_ddx}\hfill
\includegraphics[width=0.49\linewidth]{mhd_tree_dm}
\caption{Матрицы смежности и деревья: осциллятор Ван дер Поля (слева) и уравнение импульса МГД (справа). Сверху "--- истинная матрица $M$ и зашумлённая $P$ при $\alpha=0.7$, красным отмечены элементы истинной и восстановленной алгоритмом Прима матриц. Снизу "--- истинное и восстановленное деревья; zss "--- редакционное расстояние, logL "--- логарифм правдоподобия дерева по матрице $P$}
\label{fig:trees}
\end{figure}

\subsection*{Критерии точности}
\label{sec:criteria}

Критерии восстановления "--- точное совпадение матриц, совпадение наборов базисных функций после раскрытия скобок и точность описания данных. Для критерия по данным условие успеха после подбора параметров "--- $\mathrm{RMSE}/\mathrm{std}(y)<10^{-3}$. Связь критериев показана на рис.~\ref{cd:criteria}. Дополнительная мера "--- редакционное расстояние Чжана--Шаши без учёта порядка аргументов. Система считается восстановленной, если по выбранному критерию восстановлены все уравнения. Редакционное расстояние для системы "--- сумма расстояний по уравнениям.

\begin{figure}[!ht]
\centering
\begin{tikzcd}[column sep=large, row sep=large]
M \arrow[r, "\nu_\alpha"] \arrow[d, "\Phi"'] \arrow[rr, dashed, no head, bend left=35, "\text{критерий 1: } \hat M = M"] & P \arrow[r, "R"] & \hat M \arrow[d, "\Phi"]\\
f_M(\vx;\tht) \arrow[d, "\kappa"'] \arrow[dd, bend right=50, "F_D"'] & & f_{\hat M}(\vx;\tht) \arrow[d, "\kappa"] \arrow[dd, bend left=50, "F_D"]\\
\{\varphi_k\} \arrow[rr, dashed, no head, "\text{критерий 2: } ="] & & \{\hat\varphi_k\}\\
\mathbf y \arrow[rr, dashed, no head, "\text{критерий 3: } \approx"] & & f_{\hat M}(X;\hat\tht)
\end{tikzcd}
\caption{Коммутативная диаграмма эксперимента. $\nu_\alpha$ "--- зашумление с нормировкой, $R$ "--- алгоритм восстановления, $\Phi$ "--- компиляция матрицы в формулу, $\kappa$ "--- множество базисных функций, $F_D$ "--- подбор параметров и прогноз на выборке}
\label{cd:criteria}
\end{figure}

В сравнении два способа подбора параметров. В LS+LM совместно линейные параметры $\tht_L$ определяются методом наименьших квадратов при фиксированных остальных параметрах $\tht_N$. Для оптимизации по $\tht_N$ применяется метод Левенберга--Марквардта (рис.~\ref{cd:vp}). Здесь $f_0$ и $A$ "--- свободный член и матрица линейной части модели $f_0+A\tht_L$. В одноэтапном LM все параметры оптимизируются одновременно. В обоих случаях применяется мультистарт по нелинейным параметрам. Например, для частоты осциллятора $\omega$ один старт может привести к постороннему минимуму даже при верной структуре.

\begin{figure}[!ht]
\centering
\begin{tikzcd}[column sep=huge, row sep=large]
\Theta_N \arrow[r, "{(\mathrm{id},\ \hat\tht_L)}"] \arrow[dr, "\tilde S"'] & \Theta_N\times\Theta_L \arrow[d, "S"] \arrow[r, "f(X;\,\cdot\,)"] & \mathbb R^m \arrow[dl, "\|\cdot-\mathbf y\|^2"]\\
& \mathbb R_+ &
\end{tikzcd}
\caption{Метод проекции переменных: $\tilde S=S\circ(\mathrm{id},\hat\tht_L)$, $\hat\tht_L(\tht_N)=A^+(\tht_N)\bigl(\mathbf y-f_0(\tht_N)\bigr)$}
\label{cd:vp}
\end{figure}

\subsection*{Результаты}

На сводном графике показаны результаты Прима (рис.~\ref{fig:summary}). На графиках по моделям (рис.~\ref{fig:vdp}--\ref{fig:mhd}) тонкие линии соответствуют отдельным выборкам после усреднения по шуму, толстые "--- среднему по выборкам.

\begin{itemize}
\item При $\alpha\leqslant0.5$ доля точных восстановлений у DFS, BFS и алгоритма Прима составляет 100\% для всех четырёх моделей на обоих словарях. При большем шуме точность резко падает. На осцилляторе Прим опережает DFS и BFS на 1--3 п.п. Для Лоренца доли точных восстановлений у трёх алгоритмов совпадают.
\item Лишние функции снижают долю точных восстановлений. Например, у Прима при $\alpha=0.6$ для осциллятора она падает с 68 до 52\%, для Лоренца "--- с 73 до 47\%.
\item Некоторые структуры с лишними функциями всё же описывают данные. Например, вершина $\mathrm{aff}(x)=b_0+b_1x$ может воспроизвести лист $x$. Для Лотки--Вольтерры на расширенном словаре при $\alpha=0.7$ у алгоритма Прима доля точных восстановлений составляет 32\%, доля успешных аппроксимаций после подбора параметров "--- 38\%. Доли успешных аппроксимаций у LS+LM и LM в целом близки.
\item Редакционное расстояние растёт с шумом и показывает, насколько близки деревья, когда точных совпадений уже мало.
\item Усиление рёбер PCST может помочь на точном словаре. Для Лотки--Вольтерры при $\alpha=0.7$ доля точных восстановлений составляет 73\% с усилением рёбер PCST против 63\% у алгоритма Прима. На расширенном словаре PCST может включать лишние вершины. Для той же системы при $\alpha=0.5$ та же доля составляет 37\% с усилением рёбер PCST против 100\% у алгоритма Прима.
\end{itemize}

\begin{figure}[!ht]
\centering
\includegraphics[width=\linewidth]{summary}
\caption{Доля восстановленных моделей (систем) в зависимости от уровня шума, алгоритм Прима. Сплошные линии "--- точный словарь, пунктир "--- расширенный. Часть линий накладывается при совпадении долей}
\label{fig:summary}
\end{figure}

\begin{figure}[p]
\centering
\includegraphics[width=0.9\linewidth]{vdp_accuracy_exact_dict}\\[2mm]
\includegraphics[width=0.9\linewidth]{vdp_accuracy_extra_dict}
\caption{Осциллятор Ван дер Поля: точный (сверху) и расширенный (снизу) словари}
\label{fig:vdp}
\end{figure}

\begin{figure}[p]
\centering
\includegraphics[width=0.9\linewidth]{lv_accuracy_exact_dict_system}\\[2mm]
\includegraphics[width=0.9\linewidth]{lv_accuracy_extra_dict_system}
\caption{Модель Лотки--Вольтерры, система целиком: точный (сверху) и расширенный (снизу) словари. Пунктир "--- варианты на основе PCST}
\label{fig:lv}
\end{figure}

\begin{figure}[!htbp]
\centering
\includegraphics[width=0.9\linewidth]{lorenz_accuracy_exact_dict_system}\\[2mm]
\includegraphics[width=0.9\linewidth]{lorenz_accuracy_extra_dict_system}
\caption{Система Лоренца целиком, точный (сверху) и расширенный (снизу) словари. Пунктир "--- варианты на основе PCST}
\label{fig:lorenz}
\end{figure}

\begin{figure}[p]
\centering
\includegraphics[width=0.9\linewidth]{mhd_accuracy_exact_dict_system}\\[2mm]
\includegraphics[width=0.9\linewidth]{mhd_accuracy_extra_dict_system}
\caption{МГД, система целиком: точный (сверху) и расширенный (снизу) словари}
\label{fig:mhd}
\end{figure}

\end{document}
